\documentclass[12pt]{article}
\usepackage[margin=1in]{geometry}
\usepackage{amsmath}
\usepackage{amssymb}
\usepackage{graphicx}
\usepackage{tikz}
\usepackage{pgfplots}

\pgfplotsset{compat=1.17}

\title{\textbf{GED Mathematics Solutions} \\ \textbf{Slope and Equations of Lines}}
\author{\textbf{Complete Answer Key with Step-by-Step Explanations}}
\date{}

\begin{document}

\maketitle

\tableofcontents

\newpage

\section{Part A: Calculating Slope from Two Points (Questions 1-9)}

\subsection*{Question 1: Find the slope between $(2, 3)$ and $(5, 9)$}

\noindent\textbf{Solution:}

Using the slope formula: $m = \frac{y_2 - y_1}{x_2 - x_1}$

\begin{align*}
m &= \frac{9 - 3}{5 - 2} \\
m &= \frac{6}{3} \\
m &= 2
\end{align*}

\textbf{Answer: } $m = 2$

The slope is positive, indicating the line rises from left to right.

\vspace{0.3cm}

\subsection*{Question 2: Find the slope between $(-1, 4)$ and $(3, 2)$}

\noindent\textbf{Solution:}

\begin{align*}
m &= \frac{2 - 4}{3 - (-1)} \\
m &= \frac{-2}{4} \\
m &= -\frac{1}{2}
\end{align*}

\textbf{Answer: } $m = -\frac{1}{2}$

The negative slope indicates the line falls from left to right.

\vspace{0.3cm}

\subsection*{Question 3: Find the slope between $(0, 0)$ and $(4, 8)$}

\noindent\textbf{Solution:}

\begin{align*}
m &= \frac{8 - 0}{4 - 0} \\
m &= \frac{8}{4} \\
m &= 2
\end{align*}

\textbf{Answer: } $m = 2$

This line passes through the origin with a slope of 2.

\vspace{0.3cm}

\subsection*{Question 4: Find the slope between $(-3, -2)$ and $(1, 6)$}

\noindent\textbf{Solution:}

\begin{align*}
m &= \frac{6 - (-2)}{1 - (-3)} \\
m &= \frac{6 + 2}{1 + 3} \\
m &= \frac{8}{4} \\
m &= 2
\end{align*}

\textbf{Answer: } $m = 2$

\vspace{0.3cm}

\subsection*{Question 5: Find the slope between $(5, 1)$ and $(5, 7)$}

\noindent\textbf{Solution:}

\begin{align*}
m &= \frac{7 - 1}{5 - 5} \\
m &= \frac{6}{0}
\end{align*}

\textbf{Answer: } The slope is \textbf{undefined}.

Both points have the same $x$-coordinate, so the line is vertical. Vertical lines have undefined slope.

\vspace{0.3cm}

\subsection*{Question 6: Find the slope between $(2, -3)$ and $(8, -3)$}

\noindent\textbf{Solution:}

\begin{align*}
m &= \frac{-3 - (-3)}{8 - 2} \\
m &= \frac{0}{6} \\
m &= 0
\end{align*}

\textbf{Answer: } $m = 0$

Both points have the same $y$-coordinate, so the line is horizontal. Horizontal lines have zero slope.

\vspace{0.3cm}

\subsection*{Question 7: Find the slope between $(-4, 3)$ and $(2, -9)$}

\noindent\textbf{Solution:}

\begin{align*}
m &= \frac{-9 - 3}{2 - (-4)} \\
m &= \frac{-12}{6} \\
m &= -2
\end{align*}

\textbf{Answer: } $m = -2$

\vspace{0.3cm}

\subsection*{Question 8: Find the slope between $(0, 5)$ and $(-3, 2)$}

\noindent\textbf{Solution:}

\begin{align*}
m &= \frac{2 - 5}{-3 - 0} \\
m &= \frac{-3}{-3} \\
m &= 1
\end{align*}

\textbf{Answer: } $m = 1$

\vspace{0.3cm}

\subsection*{Question 9: Find the slope between $(-2, -5)$ and $(4, 7)$}

\noindent\textbf{Solution:}

\begin{align*}
m &= \frac{7 - (-5)}{4 - (-2)} \\
m &= \frac{7 + 5}{4 + 2} \\
m &= \frac{12}{6} \\
m &= 2
\end{align*}

\textbf{Answer: } $m = 2$

\newpage

\section{Part B: Identifying Parallel and Perpendicular Lines (Questions 10-17)}

\textbf{Key Concepts:}
\begin{itemize}
\item \textbf{Parallel lines} have equal slopes: $m_1 = m_2$
\item \textbf{Perpendicular lines} have slopes that are negative reciprocals: $m_1 \cdot m_2 = -1$ or $m_2 = -\frac{1}{m_1}$
\end{itemize}

\subsection*{Question 10: $m_1 = 3$ and $m_2 = 3$}

\noindent\textbf{Solution:}

Since $m_1 = m_2 = 3$, the lines are \textbf{parallel}.

\vspace{0.3cm}

\subsection*{Question 11: $m_1 = 2$ and $m_2 = -\frac{1}{2}$}

\noindent\textbf{Solution:}

Check if perpendicular: $m_1 \cdot m_2 = 2 \cdot \left(-\frac{1}{2}\right) = -1$

The lines are \textbf{perpendicular}.

\vspace{0.3cm}

\subsection*{Question 12: $m_1 = \frac{3}{4}$ and $m_2 = \frac{3}{4}$}

\noindent\textbf{Solution:}

Since $m_1 = m_2 = \frac{3}{4}$, the lines are \textbf{parallel}.

\vspace{0.3cm}

\subsection*{Question 13: $m_1 = 5$ and $m_2 = -5$}

\noindent\textbf{Solution:}

Check if perpendicular: $m_1 \cdot m_2 = 5 \cdot (-5) = -25 \neq -1$

The slopes are not equal and not negative reciprocals, so the lines are \textbf{neither parallel nor perpendicular}.

\vspace{0.3cm}

\subsection*{Question 14: $m_1 = -\frac{2}{3}$ and $m_2 = \frac{3}{2}$}

\noindent\textbf{Solution:}

Check if perpendicular: $m_1 \cdot m_2 = \left(-\frac{2}{3}\right) \cdot \frac{3}{2} = -\frac{6}{6} = -1$

The lines are \textbf{perpendicular}.

\vspace{0.3cm}

\subsection*{Question 15: $m_1 = \frac{1}{4}$ and $m_2 = -4$}

\noindent\textbf{Solution:}

Check if perpendicular: $m_1 \cdot m_2 = \frac{1}{4} \cdot (-4) = -1$

The lines are \textbf{perpendicular}.

\vspace{0.3cm}

\subsection*{Question 16: $m_1 = -1$ and $m_2 = 1$}

\noindent\textbf{Solution:}

Check if perpendicular: $m_1 \cdot m_2 = (-1) \cdot 1 = -1$

The lines are \textbf{perpendicular}.

\vspace{0.3cm}

\subsection*{Question 17: $m_1 = \frac{5}{6}$ and $m_2 = \frac{6}{5}$}

\noindent\textbf{Solution:}

Check if perpendicular: $m_1 \cdot m_2 = \frac{5}{6} \cdot \frac{6}{5} = 1 \neq -1$

The slopes are not equal and not negative reciprocals, so the lines are \textbf{neither parallel nor perpendicular}.

\newpage

\section{Part C: Finding Equations of Lines (Questions 18-27)}

\textbf{Key Form: Point-Slope Form}
$$y - y_1 = m(x - x_1)$$

This can be converted to slope-intercept form: $y = mx + b$

\subsection*{Question 18: Find the equation of a line with slope $m = 2$ passing through $(1, 3)$}

\noindent\textbf{Solution:}

Using point-slope form:
\begin{align*}
y - 3 &= 2(x - 1) \\
y - 3 &= 2x - 2 \\
y &= 2x + 1
\end{align*}

\textbf{Answer: } $y = 2x + 1$

\vspace{0.3cm}

\subsection*{Question 19: Find the equation of a line with slope $m = -3$ passing through $(2, -5)$}

\noindent\textbf{Solution:}

\begin{align*}
y - (-5) &= -3(x - 2) \\
y + 5 &= -3x + 6 \\
y &= -3x + 1
\end{align*}

\textbf{Answer: } $y = -3x + 1$

\vspace{0.3cm}

\subsection*{Question 20: Find the equation of a line passing through $(1, 2)$ and $(3, 6)$}

\noindent\textbf{Solution:}

First, find the slope:
\begin{align*}
m &= \frac{6 - 2}{3 - 1} = \frac{4}{2} = 2
\end{align*}

Using point-slope form with $(1, 2)$:
\begin{align*}
y - 2 &= 2(x - 1) \\
y - 2 &= 2x - 2 \\
y &= 2x
\end{align*}

\textbf{Answer: } $y = 2x$

\vspace{0.3cm}

\subsection*{Question 21: Find the equation of a line with slope $m = \frac{1}{2}$ passing through $(0, 4)$}

\noindent\textbf{Solution:}

Using point-slope form:
\begin{align*}
y - 4 &= \frac{1}{2}(x - 0) \\
y - 4 &= \frac{1}{2}x \\
y &= \frac{1}{2}x + 4
\end{align*}

\textbf{Answer: } $y = \frac{1}{2}x + 4$

Note: The $y$-intercept is 4 since the point $(0, 4)$ lies on the $y$-axis.

\vspace{0.3cm}

\subsection*{Question 22: Find the equation of a line passing through $(-2, 3)$ and $(4, 6)$}

\noindent\textbf{Solution:}

First, find the slope:
\begin{align*}
m &= \frac{6 - 3}{4 - (-2)} = \frac{3}{6} = \frac{1}{2}
\end{align*}

Using point-slope form with $(4, 6)$:
\begin{align*}
y - 6 &= \frac{1}{2}(x - 4) \\
y - 6 &= \frac{1}{2}x - 2 \\
y &= \frac{1}{2}x + 4
\end{align*}

\textbf{Answer: } $y = \frac{1}{2}x + 4$

\vspace{0.3cm}

\subsection*{Question 23: Find the equation of a line with slope $m = -\frac{2}{3}$ passing through $(3, 1)$}

\noindent\textbf{Solution:}

\begin{align*}
y - 1 &= -\frac{2}{3}(x - 3) \\
y - 1 &= -\frac{2}{3}x + 2 \\
y &= -\frac{2}{3}x + 3
\end{align*}

\textbf{Answer: } $y = -\frac{2}{3}x + 3$

\vspace{0.3cm}

\subsection*{Question 24: Find the equation of a line passing through $(0, 0)$ and $(5, 10)$}

\noindent\textbf{Solution:}

First, find the slope:
\begin{align*}
m &= \frac{10 - 0}{5 - 0} = \frac{10}{5} = 2
\end{align*}

Using point-slope form with $(0, 0)$:
\begin{align*}
y - 0 &= 2(x - 0) \\
y &= 2x
\end{align*}

\textbf{Answer: } $y = 2x$

This is a line through the origin.

\vspace{0.3cm}

\subsection*{Question 25: Find the equation of a line with slope $m = 0$ passing through $(2, 5)$}

\noindent\textbf{Solution:}

\begin{align*}
y - 5 &= 0(x - 2) \\
y - 5 &= 0 \\
y &= 5
\end{align*}

\textbf{Answer: } $y = 5$

This is a horizontal line at $y = 5$.

\vspace{0.3cm}

\subsection*{Question 26: Find the equation of a line passing through $(-3, -1)$ and $(2, 4)$}

\noindent\textbf{Solution:}

First, find the slope:
\begin{align*}
m &= \frac{4 - (-1)}{2 - (-3)} = \frac{5}{5} = 1
\end{align*}

Using point-slope form with $(2, 4)$:
\begin{align*}
y - 4 &= 1(x - 2) \\
y - 4 &= x - 2 \\
y &= x + 2
\end{align*}

\textbf{Answer: } $y = x + 2$

\vspace{0.3cm}

\subsection*{Question 27: Find the equation of a line with slope $m = -1$ passing through $(4, -2)$}

\noindent\textbf{Solution:}

\begin{align*}
y - (-2) &= -1(x - 4) \\
y + 2 &= -x + 4 \\
y &= -x + 2
\end{align*}

\textbf{Answer: } $y = -x + 2$

\newpage

\section{Part D: Graphing Linear Equations (Questions 28-33)}

\textbf{Key Concept: Slope-Intercept Form}
$$y = mx + b$$

where $m$ is the slope and $b$ is the $y$-intercept.

\subsection*{Question 28: $2x + y = 5$}

\noindent\textbf{Solution:}

Solve for $y$:
\begin{align*}
2x + y &= 5 \\
y &= -2x + 5
\end{align*}

\textbf{Slope: } $m = -2$

\textbf{$y$-intercept: } $b = 5$ (point: $(0, 5)$)

\vspace{0.3cm}

\subsection*{Question 29: $3x - y = 6$}

\noindent\textbf{Solution:}

Solve for $y$:
\begin{align*}
3x - y &= 6 \\
-y &= -3x + 6 \\
y &= 3x - 6
\end{align*}

\textbf{Slope: } $m = 3$

\textbf{$y$-intercept: } $b = -6$ (point: $(0, -6)$)

\vspace{0.3cm}

\subsection*{Question 30: $4x + 2y = 8$}

\noindent\textbf{Solution:}

Solve for $y$:
\begin{align*}
4x + 2y &= 8 \\
2y &= -4x + 8 \\
y &= -2x + 4
\end{align*}

\textbf{Slope: } $m = -2$

\textbf{$y$-intercept: } $b = 4$ (point: $(0, 4)$)

\vspace{0.3cm}

\subsection*{Question 31: $y = 3x - 2$}

\noindent\textbf{Solution:}

This equation is already in slope-intercept form.

\textbf{Slope: } $m = 3$

\textbf{$y$-intercept: } $b = -2$ (point: $(0, -2)$)

\vspace{0.3cm}

\subsection*{Question 32: $-2x + 3y = 9$}

\noindent\textbf{Solution:}

Solve for $y$:
\begin{align*}
-2x + 3y &= 9 \\
3y &= 2x + 9 \\
y &= \frac{2}{3}x + 3
\end{align*}

\textbf{Slope: } $m = \frac{2}{3}$

\textbf{$y$-intercept: } $b = 3$ (point: $(0, 3)$)

\vspace{0.3cm}

\subsection*{Question 33: $5x - 5y = 10$}

\noindent\textbf{Solution:}

Solve for $y$:
\begin{align*}
5x - 5y &= 10 \\
-5y &= -5x + 10 \\
y &= x - 2
\end{align*}

\textbf{Slope: } $m = 1$

\textbf{$y$-intercept: } $b = -2$ (point: $(0, -2)$)

\newpage

\section{Part E: Real-World Applications of Slope (Questions 34-38)}

\subsection*{Question 34: A wheelchair ramp rises 2 feet over a horizontal distance of 24 feet. What is the slope of the ramp?}

\noindent\textbf{Solution:}

Slope is defined as the ratio of rise (vertical change) to run (horizontal change).

\begin{align*}
\text{Slope} &= \frac{\text{rise}}{\text{run}} \\
&= \frac{2}{24} \\
&= \frac{1}{12}
\end{align*}

\textbf{Answer: } $m = \frac{1}{12} \approx 0.083$ or about 8.3\%

The slope is very gentle, which is appropriate for wheelchair accessibility (ADA standard is typically 1:12).

\vspace{0.3cm}

\subsection*{Question 35: A staircase has steps that rise 7 inches for every 10 inches of horizontal distance. What is the slope of the staircase?}

\noindent\textbf{Solution:}

\begin{align*}
\text{Slope} &= \frac{\text{rise}}{\text{run}} \\
&= \frac{7}{10} \\
&= 0.7
\end{align*}

\textbf{Answer: } $m = \frac{7}{10}$ or $m = 0.7$

\vspace{0.3cm}

\subsection*{Question 36: A road elevation increases 300 feet for every 2000 feet of horizontal distance. Express this as a slope in simplest form.}

\noindent\textbf{Solution:}

\begin{align*}
\text{Slope} &= \frac{300}{2000}
\end{align*}

Simplify by finding the GCD of 300 and 2000:
\begin{align*}
\frac{300}{2000} &= \frac{3}{20}
\end{align*}

\textbf{Answer: } $m = \frac{3}{20}$ or $m = 0.15$ or 15\%

\vspace{0.3cm}

\subsection*{Question 37: A mountain trail descends 400 feet in elevation over a horizontal distance of 1 mile (5280 feet). What is the slope of the trail?}

\noindent\textbf{Solution:}

Since the trail descends (goes down), the slope is negative.

\begin{align*}
\text{Slope} &= -\frac{400}{5280}
\end{align*}

Simplify by finding the GCD of 400 and 5280:
\begin{align*}
-\frac{400}{5280} &= -\frac{5}{66}
\end{align*}

\textbf{Answer: } $m = -\frac{5}{66} \approx -0.0758$ or about $-7.58\%$

\vspace{0.3cm}

\subsection*{Question 38: A ladder leans against a wall, reaching 12 feet high. The base of the ladder is 5 feet from the wall. What is the slope of the ladder?}

\noindent\textbf{Solution:}

The ladder reaches from the point $(0, 0)$ (base on ground) to the point $(5, 12)$ (top of wall).

\begin{align*}
\text{Slope} &= \frac{12 - 0}{5 - 0} \\
&= \frac{12}{5} \\
&= 2.4
\end{align*}

\textbf{Answer: } $m = \frac{12}{5}$ or $m = 2.4$

The slope is steep because the ladder rises much more than it extends horizontally.

\newpage

\section{Part F: Distance Between Points (Questions 39-43)}

\textbf{Distance Formula:}
$$d = \sqrt{(x_2 - x_1)^2 + (y_2 - y_1)^2}$$

\subsection*{Question 39: Find the distance between $(0, 0)$ and $(3, 4)$}

\noindent\textbf{Solution:}

\begin{align*}
d &= \sqrt{(3 - 0)^2 + (4 - 0)^2} \\
&= \sqrt{3^2 + 4^2} \\
&= \sqrt{9 + 16} \\
&= \sqrt{25} \\
&= 5
\end{align*}

\textbf{Answer: } $d = 5$ units

This is the famous 3-4-5 Pythagorean triple.

\vspace{0.3cm}

\subsection*{Question 40: Find the distance between $(1, 2)$ and $(4, 6)$}

\noindent\textbf{Solution:}

\begin{align*}
d &= \sqrt{(4 - 1)^2 + (6 - 2)^2} \\
&= \sqrt{3^2 + 4^2} \\
&= \sqrt{9 + 16} \\
&= \sqrt{25} \\
&= 5
\end{align*}

\textbf{Answer: } $d = 5$ units

\vspace{0.3cm}

\subsection*{Question 41: Find the distance between $(-2, 1)$ and $(3, 5)$}

\noindent\textbf{Solution:}

\begin{align*}
d &= \sqrt{(3 - (-2))^2 + (5 - 1)^2} \\
&= \sqrt{(5)^2 + (4)^2} \\
&= \sqrt{25 + 16} \\
&= \sqrt{41}
\end{align*}

\textbf{Answer: } $d = \sqrt{41} \approx 6.4$ units

\vspace{0.3cm}

\subsection*{Question 42: Find the distance between $(0, -3)$ and $(4, 0)$}

\noindent\textbf{Solution:}

\begin{align*}
d &= \sqrt{(4 - 0)^2 + (0 - (-3))^2} \\
&= \sqrt{4^2 + 3^2} \\
&= \sqrt{16 + 9} \\
&= \sqrt{25} \\
&= 5
\end{align*}

\textbf{Answer: } $d = 5$ units

Another 3-4-5 Pythagorean triple (3, 4, 5).

\vspace{0.3cm}

\subsection*{Question 43: Find the distance between $(-1, -1)$ and $(2, 3)$}

\noindent\textbf{Solution:}

\begin{align*}
d &= \sqrt{(2 - (-1))^2 + (3 - (-1))^2} \\
&= \sqrt{(3)^2 + (4)^2} \\
&= \sqrt{9 + 16} \\
&= \sqrt{25} \\
&= 5
\end{align*}

\textbf{Answer: } $d = 5$ units

Yet another 3-4-5 Pythagorean triple.

\newpage

\section{Part G: Systems of Equations (Questions 44-48)}

\textbf{Methods: Substitution or Elimination}

\subsection*{Question 44: $\begin{cases} y = 2x + 1 \\ y = -x + 4 \end{cases}$}

\noindent\textbf{Solution (Substitution Method):}

Since both expressions equal $y$, set them equal to each other:
\begin{align*}
2x + 1 &= -x + 4 \\
3x + 1 &= 4 \\
3x &= 3 \\
x &= 1
\end{align*}

Substitute $x = 1$ into the first equation:
\begin{align*}
y &= 2(1) + 1 \\
y &= 3
\end{align*}

\textbf{Answer: } $(x, y) = (1, 3)$

\vspace{0.3cm}

\subsection*{Question 45: $\begin{cases} x + y = 5 \\ x - y = 1 \end{cases}$}

\noindent\textbf{Solution (Elimination Method):}

Add the two equations to eliminate $y$:
\begin{align*}
(x + y) + (x - y) &= 5 + 1 \\
2x &= 6 \\
x &= 3
\end{align*}

Substitute $x = 3$ into the first equation:
\begin{align*}
3 + y &= 5 \\
y &= 2
\end{align*}

\textbf{Answer: } $(x, y) = (3, 2)$

\vspace{0.3cm}

\subsection*{Question 46: $\begin{cases} 2x + y = 7 \\ x - y = 2 \end{cases}$}

\noindent\textbf{Solution (Elimination Method):}

Add the two equations to eliminate $y$:
\begin{align*}
(2x + y) + (x - y) &= 7 + 2 \\
3x &= 9 \\
x &= 3
\end{align*}

Substitute $x = 3$ into the second equation:
\begin{align*}
3 - y &= 2 \\
-y &= -1 \\
y &= 1
\end{align*}

\textbf{Answer: } $(x, y) = (3, 1)$

\vspace{0.3cm}

\subsection*{Question 47: $\begin{cases} 3x - 2y = 4 \\ x + y = 5 \end{cases}$}

\noindent\textbf{Solution (Substitution Method):}

From the second equation, solve for $y$:
\begin{align*}
x + y &= 5 \\
y &= 5 - x
\end{align*}

Substitute into the first equation:
\begin{align*}
3x - 2(5 - x) &= 4 \\
3x - 10 + 2x &= 4 \\
5x - 10 &= 4 \\
5x &= 14 \\
x &= \frac{14}{5}
\end{align*}

Substitute back to find $y$:
\begin{align*}
y &= 5 - \frac{14}{5} \\
y &= \frac{25}{5} - \frac{14}{5} \\
y &= \frac{11}{5}
\end{align*}

\textbf{Answer: } $(x, y) = \left(\frac{14}{5}, \frac{11}{5}\right)$ or $(2.8, 2.2)$

\vspace{0.3cm}

\subsection*{Question 48: $\begin{cases} y = \frac{1}{2}x + 3 \\ y = -2x + 8 \end{cases}$}

\noindent\textbf{Solution (Substitution Method):}

Set the two expressions for $y$ equal:
\begin{align*}
\frac{1}{2}x + 3 &= -2x + 8 \\
\frac{1}{2}x + 2x &= 8 - 3 \\
\frac{1}{2}x + \frac{4}{2}x &= 5 \\
\frac{5}{2}x &= 5 \\
x &= 2
\end{align*}

Substitute $x = 2$ into the first equation:
\begin{align*}
y &= \frac{1}{2}(2) + 3 \\
y &= 1 + 3 \\
y &= 4
\end{align*}

\textbf{Answer: } $(x, y) = (2, 4)$

\newpage

\section{Part H: Mixed Review (Questions 49-50)}

\subsection*{Question 49: Two parallel lines have equations $y = 3x + 2$ and $y = 3x - 5$. Find the perpendicular slope to these lines, and write an equation of a line that passes through $(2, 4)$ with that perpendicular slope.}

\noindent\textbf{Solution:}

The parallel lines have slope $m = 3$.

For a perpendicular line, the slope is the negative reciprocal:
\begin{align*}
m_{\perp} &= -\frac{1}{3}
\end{align*}

Now write the equation of a line with slope $m = -\frac{1}{3}$ passing through $(2, 4)$:
\begin{align*}
y - 4 &= -\frac{1}{3}(x - 2) \\
y - 4 &= -\frac{1}{3}x + \frac{2}{3} \\
y &= -\frac{1}{3}x + \frac{2}{3} + 4 \\
y &= -\frac{1}{3}x + \frac{2}{3} + \frac{12}{3} \\
y &= -\frac{1}{3}x + \frac{14}{3}
\end{align*}

\textbf{Answer:}
\begin{itemize}
\item Perpendicular slope: $m = -\frac{1}{3}$
\item Equation: $y = -\frac{1}{3}x + \frac{14}{3}$
\end{itemize}

\vspace{0.5cm}

\subsection*{Question 50: A line passes through the points $(-2, 1)$ and $(4, 7)$. Find the slope, write the equation in slope-intercept form, and find the distance between the two points.}

\noindent\textbf{Solution:}

\textbf{Part 1: Find the slope}
\begin{align*}
m &= \frac{7 - 1}{4 - (-2)} \\
&= \frac{6}{6} \\
&= 1
\end{align*}

\textbf{Part 2: Write the equation in slope-intercept form}

Using point-slope form with point $(4, 7)$:
\begin{align*}
y - 7 &= 1(x - 4) \\
y - 7 &= x - 4 \\
y &= x + 3
\end{align*}

\textbf{Part 3: Find the distance between the two points}
\begin{align*}
d &= \sqrt{(4 - (-2))^2 + (7 - 1)^2} \\
&= \sqrt{(6)^2 + (6)^2} \\
&= \sqrt{36 + 36} \\
&= \sqrt{72} \\
&= \sqrt{36 \cdot 2} \\
&= 6\sqrt{2}
\end{align*}

\textbf{Answer:}
\begin{itemize}
\item Slope: $m = 1$
\item Equation: $y = x + 3$
\item Distance: $d = 6\sqrt{2} \approx 8.49$ units
\end{itemize}

\end{document}